% !TEX program = lualatex
% Auto-generated from syllabus — 09: 情報の加工2：表計算とデータ整理

\documentclass[aspectratio=43,professionalfonts,handout,10pt]{beamer}
\usepackage{beamer_template}

\newcommand{\LectureNum}{09}
\newcommand{\LectureTitle}{情報の加工2：表計算とデータ整理}
\newcommand{\CourseName}{情報リテラシー}
\newcommand{\TextPages}{-}

\begin{document}
% --- Slide 1: title ---

\begin{frame}{\CourseName\quad 第\LectureNum 回「\LectureTitle 」}
  {\small \TermName\quad \TeacherName\quad ／\quad 教科書 \TextPages}

  \vfill

  \begin{block}{今日の流れ}
    \begin{enumerate}
      \item データの取得と操作
      \item 記述統計：Excelの統計関数
      \item 統計量の意味と計算式
      \item グラフによる可視化
    \end{enumerate}
  \end{block}
  \vfill

  \begin{block}{今日の目標}
    \begin{itemize}
      \item 表計算ソフトを用いてデータの整理・集計・グラフによる可視化ができる。
    \end{itemize}
  \end{block}
  \vfill

  \begin{exampleblock}{キーワード}
    \small \textbf{Excel} ／ \textbf{CSV} ／ \textbf{オープンデータ} ／ \textbf{統計関数}
  \end{exampleblock}
\end{frame}

% --- Slide 2: データの取得と操作 ---

\begin{frame}{データの取得と操作}
  \begin{exampleblock}{最初にやること：保存先の設定}
    「ファイル」→「オプション」→「保存」→\textbf{「既定でコンピューターに保存する」}にチェック
  \end{exampleblock}

  \begin{block}{データ取得の方法}
    \begin{itemize}
      \item 手動入力：直接セルに入力する基本的な方法
      \item CSVファイル読み込み：「データ」→「テキストまたはCSVから」
      \item オープンデータ：e-Stat（政府統計）・気象庁・国土地理院等を活用
    \end{itemize}
  \end{block}

  \begin{slidenote}
    CSVはComma-Separated Valuesの略。テキストファイルなのでExcel以外のソフトでも開ける。
    オープンデータは自由に利用・加工・再配布が可能なデータ。e-Statでは人口・経済・労働などの統計データが取得できる。
    データ整形の基本：不要な行・列の削除，並べ替え（「データ」タブ→「昇順」/「降順」）。
  \end{slidenote}
\end{frame}

% --- Slide 3: 記述統計 ---

\begin{frame}{記述統計：Excelの統計関数}
  \begin{block}{よく使う統計関数}
    \footnotesize
    \begin{columns}[T]
      \begin{column}{0.48\textwidth}
        \begin{itemize}
          \item 合計：\texttt{=SUM(B2:B11)}
          \item 個数：\texttt{=COUNT(B2:B11)}
          \item 四捨五入：\texttt{=ROUND(B2,1)}
          \item 平均値：\texttt{=AVERAGE(B2:B11)}
        \end{itemize}
      \end{column}
      \begin{column}{0.48\textwidth}
        \begin{itemize}
          \item 中央値：\texttt{=MEDIAN(B2:B11)}
          \item 最頻値：\texttt{=MODE.SNGL(B2:B11)}
          \item 標準偏差：\texttt{=STDEV.S(B2:B11)}
          \item 最大・最小：\texttt{=MAX()} / \texttt{=MIN()}
        \end{itemize}
      \end{column}
    \end{columns}
  \end{block}

  \begin{slidenote}
    STDEV.S の S は Sample（標本）。母集団全体なら STDEV.P。
    ROUND(B2,1)→小数第1位，ROUND(B2,0)→整数に四捨五入。
  \end{slidenote}
\end{frame}

% --- Slide 4: 統計量の意味と計算式 ---

\begin{frame}{統計量の意味と計算式}
  \footnotesize
  \begin{tabular}{lll}
    \hline
    統計量 & 意味 & 計算式 \\
    \hline
    合計 & 全データの和 & $\displaystyle\sum_{i=1}^{n} x_i$ \\[4pt]
    個数 & データの件数 & $n$ \\[2pt]
    平均値 & データの重心。外れ値の影響を受けやすい & $\displaystyle\bar{x} = \frac{1}{n}\sum_{i=1}^{n} x_i$ \\[4pt]
    中央値 & 昇順に並べた真ん中の値。外れ値に強い & 順位 $\frac{n+1}{2}$ 番目の値 \\[4pt]
    最頻値 & 最も多く現れる値 & 出現回数が最大の値 \\[2pt]
    最大・最小 & データの範囲確認 & $\max(x_i)$，$\min(x_i)$ \\[2pt]
    標準偏差 & ばらつきの大きさ。大きいほど散らばっている & $\displaystyle s = \sqrt{\frac{1}{n-1}\sum_{i=1}^{n}(x_i - \bar{x})^2}$ \\[4pt]
    \hline
  \end{tabular}

  \begin{slidenote}
    式は「こういう計算をしている」という雰囲気をつかむだけでよい。Excelが自動で計算してくれる。
    標準偏差の式：各データが平均からどれだけ離れているかの「平均的な距離」。
  \end{slidenote}
\end{frame}

% --- Slide 5: グラフによる可視化 ---

\begin{frame}{グラフによる可視化}
  \begin{block}{グラフ作成の手順}
    \begin{enumerate}
      \item グラフにしたいデータ範囲を選択する
      \item 「挿入」タブ→「グラフ」から種類を選ぶ
      \item タイトル・軸ラベルを設定する
    \end{enumerate}
  \end{block}

  \begin{exampleblock}{グラフの使い分け}
    \begin{itemize}
      \item \textbf{棒グラフ}：項目ごとの量を比較したいとき
      \item \textbf{折れ線グラフ}：時間の変化を見たいとき
      \item \textbf{円グラフ}：全体に対する割合を見たいとき
      \item \textbf{散布図}：2つの数値の関係を見たいとき
    \end{itemize}
  \end{exampleblock}

  \begin{slidenote}
    グラフはデータを「見える化」するツール。同じデータでも種類を変えると印象が変わる。
    軸ラベル・単位・タイトルがないグラフは読み手に伝わらない。必ず設定する習慣をつけよう。
    散布図：2つの数値の関係（相関）を見るグラフ。右上がりなら正の相関。
  \end{slidenote}
\end{frame}

% --- Slide 6: 演習（提出課題）---

\begin{frame}{演習（提出課題）}
  \begin{alertblock}{課題：Teamsからファイルをダウンロードして完成させ，提出せよ}
    \begin{enumerate}\itemsep=2pt
      \item 数学（B列）の統計量を求める\\
            （合計・個数・平均値・中央値・最頻値・最大値・最小値・標準偏差）
      \item 平均値・標準偏差は小数第2位まで四捨五入する
      \item 数学と英語の点数で散布図を作成し，タイトル・軸ラベルを設定する
    \end{enumerate}
  \end{alertblock}

  \begin{exampleblock}{ファイル名}
    \texttt{1E\_名前\_演習2.xlsx}（例：\texttt{1E\_山田太郎\_演習2.xlsx}）
  \end{exampleblock}

  \begin{slidenote}
    提出前にローカル保存確認。散布図はシート内に埋め込んだままでよい。
  \end{slidenote}
\end{frame}

% --- チェックリスト ---

\begin{frame}{まとめ・チェックリスト}
  \begin{block}{自己チェック（できたら $\checkmark$ をつけよう）}
    \begin{itemize}
      \item[$\square$] Excelでデータ入力・並べ替えができた
      \item[$\square$] SUM・COUNT・ROUND・AVERAGE・MEDIAN・STDEV.S等を使えた
      \item[$\square$] 散布図を作成し，タイトル・軸ラベルを設定できた
      \item[$\square$] 演習をTeamsの課題から提出できた
    \end{itemize}
  \end{block}

  \vfill

  {\small 演習のExcelファイルをTeamsの課題から提出すること。}
\end{frame}

\end{document}
